\documentclass[12pt,a4paper]{article}
\usepackage{amsmath,amssymb,amsthm}
\usepackage[colorlinks=true,linkcolor=blue,citecolor=blue,urlcolor=blue]{hyperref}
\usepackage{geometry}
\geometry{margin=2.5cm}
\usepackage{microtype}

% -------------------------------------------------------
\title{\textbf{BEC Sonic Horizon as a First Experimental Falsifier\\
for the ECI Krylov-Complexity / Type~II$_\infty$ Framework}\\[0.5em]
\large Proposal Note, ECI v6.0.20 --- DRAFT}
\author{ECI Collaboration internal note\\
(F1 agent, 2026-05-03)}
\date{\today}
% -------------------------------------------------------

\newtheorem{conjecture}{Conjecture}
\newtheorem{remark}{Remark}
\renewcommand{\TH}{T_{\mathrm{H}}}
\newcommand{\lambdaL}{\lambda_L^{\mathrm{sonic}}}
\newcommand{\kB}{k_{\mathrm{B}}}

\begin{document}
\maketitle

\begin{abstract}
We propose using the Bose--Einstein condensate (BEC) sonic black hole of
Steinhauer et al.\ as the first experimental testbed for the ECI
Krylov-complexity / type~II$_\infty$ framework.
The BEC phonon algebra on the exterior (subsonic) region is generically
type~III$_1$ in the thermodynamic limit; a crossed product with the
modular flow promotes it to type~II$_\infty$, following Witten
\cite{witten2022}.
The modular Lyapunov exponent saturates at $\lambdaL = 2\pi \kB \TH /
\hbar$ (Caputa--Mag\'an--Patramanis--Tonni \cite{caputa2023}), which
equals the acoustic surface gravity $\kappa$.
For Steinhauer's 2019 system \cite{steinhauer2019}, we compute
$\lambdaL \approx 2.88 \times 10^2\,\mathrm{s}^{-1}$ and a
Krylov saturation time $t_K \approx 3.5\,\mathrm{ms}$.
The non-trivial ECI prediction is that the temporal decay rate of the
phonon $g^{(2)}(\tau)$ function equals $\lambdaL$ independently of the
thermality measurement.
This has not been measured.
We also flag: (i)~the ECI~v6.0.10 ``8.29\%~saturation'' benchmark is
\emph{unverifiable} and should be retracted; (ii)~no Steinhauer-group
2024--2026 saturation-envelope paper exists on arXiv; (iii)~the
Krylov-Diameter Theorem conjecture gives only an order-of-magnitude
estimate.
\end{abstract}

\tableofcontents
\bigskip

%=======================================================
\section{BEC Sonic Horizon: Experimental Setup}
\label{sec:setup}
%=======================================================

\subsection{Steinhauer's analogue black hole}

Steinhauer and collaborators at the Technion have constructed a series of
sonic black holes in atomic $^{87}$Rb Bose--Einstein condensates
\cite{steinhauer2016,steinhauer2019,kolobov2021}.
The key construction: a trapped, elongated BEC is accelerated through a
step potential so that the condensate flow velocity $v(x)$ exceeds the
local speed of sound $c_s(x)$ beyond a critical point $x_H$.
The locus $v(x_H) = c_s(x_H)$ defines the sonic horizon.

Phonons (quantized sound waves) propagating in the supersonic interior
region cannot escape against the background flow; they are the analogue of
photons inside an event horizon.
The acoustic metric,
\begin{equation}
  g_{\mu\nu}^{\mathrm{ac}} dx^\mu dx^\nu
  = \frac{\rho_0}{c_s}
    \Bigl[-(c_s^2 - v^2)\,dt^2 - 2v\,dt\,dx + dx^2 + dy^2 + dz^2\Bigr],
\end{equation}
has a horizon structure precisely analogous to the Schwarzschild metric
near $r = 2GM$ \cite{unruh1981,delhom2025}.

\subsection{Measured Hawking temperature}

The analogue surface gravity is
\begin{equation}
  \kappa = \tfrac{1}{2}\bigl|\partial_x(v - c_s)\bigr|_{x=x_H}.
\end{equation}
Hawking's prediction gives an analogue temperature
\begin{equation}
  \TH = \frac{\hbar\,\kappa}{2\pi\,\kB}.
\end{equation}
Steinhauer~et~al.\ (2019, arXiv:1809.00913) measured the correlation
spectrum of phonon Hawking pairs and found it to be thermal at temperature
\begin{equation}
  \TH = (0.35 \pm 0.10)\,\mathrm{nK},
  \label{eq:TH_measured}
\end{equation}
in quantitative agreement with the acoustic surface gravity implied by the
density and velocity profiles.
The Hawking radiation is in the linear-dispersion (relativistic) regime,
confirming the analogy with a real black hole.

\subsection{Phonon field and UV cutoff}

The BEC phonon field $\phi$ satisfies the Bogoliubov--de Gennes (BdG)
Hamiltonian:
\begin{equation}
  \hat{H}_{\mathrm{BdG}} = \int \frac{dk}{2\pi}
    \Bigl[\omega_k\,\hat{a}_k^\dagger \hat{a}_k
    + \text{(pair terms)}\Bigr],
  \quad
  \omega_k^2 = c_s^2 k^2 + \Bigl(\frac{\hbar k^2}{2m}\Bigr)^2.
\end{equation}
The healing length $\xi = \hbar/\sqrt{2m\mu}$ (where $\mu$ is the chemical
potential) sets the UV cutoff: for $k\xi \ll 1$, the dispersion is
linear (phononic) and the acoustic metric description holds; for
$k\xi \sim 1$, the dispersion becomes quadratic.
Typical parameters: $c_s \sim 0.5\,\mathrm{mm/s}$, $\xi \sim 0.2$--$1\,\mu\mathrm{m}$.

The BEC naturally resolves the trans-Planckian problem \cite{jacobson1991}:
only phonons with $k \ll 1/\xi$ contribute to Hawking radiation, so no
trans-Planckian modes are involved.

%=======================================================
\section{ECI Apparatus: Type~II$_\infty$ for the BEC Observer Subregion}
\label{sec:algebra}
%=======================================================

\subsection{Type~III$_1$ algebra of BEC phonons at the horizon}

Consider the algebra $\mathcal{A}_{\mathrm{ext}}$ of phonon observables
localised in the subsonic (exterior) region $x < x_H$.
In the thermodynamic limit ($L \to \infty$, $\xi \to 0$), this algebra is
of type~III$_1$ in the Tomita--Takesaki classification, for the following
reasons:
\begin{enumerate}
  \item The Bogoliubov vacuum $|\Omega_{\mathrm{BdG}}\rangle$ is a KMS
    state on $\mathcal{A}_{\mathrm{ext}}$ at inverse temperature
    $\beta_H = 1/\kB \TH$ (Hawking temperature), analogous to the
    Bisognano--Wichmann theorem for Rindler wedge algebras.
  \item The modular operator $\Delta_\Omega$ has continuous spectrum
    $\mathbb{R}_{>0}$ without eigenvalues in the thermodynamic limit.
  \item The flow $t \mapsto \Delta_\Omega^{it}$ implements the acoustic
    Rindler boost (Killing flow generated by $\partial_t + \kappa x
    \partial_x$ near the horizon), which is unbounded.
\end{enumerate}
The type~III$_1$ classification is the generic QFT result, also obtained
for the Rindler wedge and Schwarzschild black holes.

\subsection{Crossed product and type~II$_\infty$}

Following Witten \cite{witten2022}, one forms the crossed product
\begin{equation}
  \widetilde{\mathcal{A}}_{\mathrm{ext}}
    = \mathcal{A}_{\mathrm{ext}} \rtimes_{\sigma^\Omega} \mathbb{R},
\end{equation}
where $\sigma^\Omega_t = \Delta_\Omega^{it}(\,\cdot\,)\Delta_\Omega^{-it}$
is the modular automorphism group.
The crossed product $\widetilde{\mathcal{A}}_{\mathrm{ext}}$ is of
type~II$_\infty$: it has a well-defined (but state-dependent) trace, and
entropy differences between states are finite and state-independent.

In the ECI context \cite{eci_krylov_diameter}, the additional ingredient
is the modular Hamiltonian
\begin{equation}
  K_{\mathrm{BEC}} = -\log\Delta_\Omega
    = -\log\rho_{\mathrm{ext}},
\end{equation}
where $\rho_{\mathrm{ext}} = \mathrm{Tr}_{\mathrm{int}}|\Omega\rangle\langle\Omega|$
is the reduced density matrix on the exterior.
Explicitly, in the Unruh decomposition:
\begin{equation}
  K_{\mathrm{BEC}} = \frac{2\pi}{\kappa}
    \int_0^\infty dk\,\omega_k\,\hat{a}_k^\dagger \hat{a}_k + \text{const.}
\end{equation}

\subsection{Finite-$\xi$ correction}

At finite healing length $\xi > 0$, the algebra $\mathcal{A}_{\mathrm{ext}}$
is not exactly type~III$_1$: the UV cutoff renders the operator algebra
effectively type~I on scales $k > 1/\xi$.
The type~II$_\infty$ structure of $\widetilde{\mathcal{A}}_{\mathrm{ext}}$
is therefore an \emph{infrared} property, valid for phonon modes with
$k \xi \ll 1$.
This is consistent with the Hawking radiation being in the linear-dispersion
regime.

\begin{remark}
The BEC quasi-vacuum $|\Omega_{\mathrm{BdG}}\rangle$ is \emph{not} in the
same Verch folium as the conformal vacuum used in the ECI $K_{\mathrm{FRW}}$
carve-out construction.
The Bogoliubov transformation that relates the two vacua is non-trivial at
the healing length scale.
However, for phonon modes $k\xi \ll 1$, the two vacua share the same
KMS structure at $\TH$, which suffices for the Krylov-complexity analysis.
\end{remark}

%=======================================================
\section{Krylov-Complexity Prediction at the Sonic Horizon}
\label{sec:krylov}
%=======================================================

\subsection{Modular Lyapunov saturation}

Caputa, Mag\'an, Patramanis and Tonni \cite{caputa2023} prove that for an
operator $\mathcal{O}$ evolved under modular flow in a type~III$_1$
algebra, the spread complexity (Krylov complexity $C_K$) grows as:
\begin{equation}
  C_K(t_{\mathrm{mod}}) \sim e^{\lambda_L^{\mathrm{mod}}\,t_{\mathrm{mod}}},
  \qquad \lambda_L^{\mathrm{mod}} = 2\pi,
\end{equation}
at late modular time $t_{\mathrm{mod}} \gg 1$.
Here, $\lambda_L^{\mathrm{mod}} = 2\pi$ is the universal modular Lyapunov
exponent in units where the local temperature is unity.

Converting to laboratory time via $t_{\mathrm{lab}} =
t_{\mathrm{mod}}/\kappa$, the lab-frame Lyapunov exponent is:
\begin{equation}
  \boxed{\lambdaL = \kappa = \frac{2\pi\,\kB\,\TH}{\hbar}.}
  \label{eq:lambdaL}
\end{equation}
This saturates the Maldacena--Shenker--Stanford (MSS) chaos bound
$\lambda_L \leq 2\pi\kB T/\hbar$ \cite{mss2016}, consistent with the
horizon being a maximally scrambling surface.

\subsection{Krylov-Diameter Theorem conjecture}

The ECI Krylov-Diameter Theorem (Theorem~4) predicts
\begin{equation}
  \frac{dC_K}{dt}\Biggl|_{\mathrm{horizon}}
    = \frac{1}{R_{\mathrm{proper}}(\eta_c)},
\end{equation}
where $R_{\mathrm{proper}}$ is the proper size of the past-light-cone
diamond at the horizon.
In the BEC analogue, the natural identification is
$R_{\mathrm{proper}} \sim \xi$ (healing length), giving
\begin{equation}
  \lambdaL^{\mathrm{KD}} \sim \frac{c_s}{\xi}.
  \label{eq:KD_conjecture}
\end{equation}

\begin{remark}[Honest assessment of \eqref{eq:KD_conjecture}]
The numerical computation (see \texttt{sympy\_check.py}) gives
$c_s/\xi \approx 10^3\,\mathrm{s}^{-1}$ for $c_s = 0.5\,\mathrm{mm/s}$,
$\xi = 0.5\,\mu\mathrm{m}$, corresponding to $\TH \approx 1.2\,\mathrm{nK}$.
Steinhauer measured $\TH = 0.35\,\mathrm{nK}$, a factor of $\sim 3.5$
discrepancy.
The KD conjecture is therefore order-of-magnitude only; the actual surface
gravity depends on the flow-profile geometry at the horizon, not just $\xi$.
This conjecture \emph{cannot} serve as a precision test until the
flow-profile derivative $\partial_x(v - c_s)|_{x_H}$ is specified.
\end{remark}

\subsection{Numerical results for Steinhauer's 2019 parameters}

Using $\TH = 0.35\,\mathrm{nK}$ from \cite{steinhauer2019}:
\begin{align}
  \kappa &= \frac{2\pi\,\kB\,\TH}{\hbar}
    \approx 2.88 \times 10^2\,\mathrm{rad\,s}^{-1},\\
  \lambdaL &= \kappa \approx 2.88 \times 10^2\,\mathrm{s}^{-1},\\
  t_K &= \frac{1}{\lambdaL} \approx 3.5\,\mathrm{ms}
    \quad [3.5 \times 10^3\,\mu\mathrm{s}],\\
  \nu_{\mathrm{char}} &= \frac{\kB\,\TH}{h} \approx 7.3\,\mathrm{Hz}.
\end{align}

The Krylov saturation time $t_K \approx 3.5\,\mathrm{ms}$ is well within
the coherence lifetime of a $^{87}$Rb BEC (typically $10$--$100\,\mathrm{ms}$),
suggesting the saturation is in principle observable.

%=======================================================
\section{Experimental Signature: Phonon Coherence Decay Rate}
\label{sec:observable}
%=======================================================

\subsection{ECI prediction: $g^{(2)}(\tau)$ decay rate}

The central ECI prediction for the BEC sonic horizon is:

\begin{conjecture}[Modular Lyapunov in $g^{(2)}(\tau)$]
The second-order coherence function of phonons near the sonic horizon,
\begin{equation}
  g^{(2)}(\tau) = \frac{\langle\hat{n}(t)\hat{n}(t+\tau)\rangle}
    {\langle\hat{n}(t)\rangle^2},
\end{equation}
decays at late times as
\begin{equation}
  g^{(2)}(\tau) - 1 \propto e^{-\Gamma\,\tau},
  \qquad
  \Gamma = \lambdaL = \frac{2\pi\,\kB\,\TH}{\hbar},
  \label{eq:prediction}
\end{equation}
where $\TH$ is the Hawking temperature measured \emph{independently}
from the correlation spectrum.
\end{conjecture}

This is a non-trivial prediction because it specifies the temporal decay
\emph{rate}, not just the thermality of the spectrum.
Standard Hawking-radiation observations (Steinhauer 2019) establish thermality
at $\TH$ from the frequency distribution; they do not measure the temporal
decay of $g^{(2)}(\tau)$ against this same $\TH$.

\subsection{Phonon spectral measure}

The ECI framework predicts the phonon spectral measure
\begin{equation}
  \mu_\psi(\omega) \sim R(\omega)\,e^{-2\pi\omega/\kappa},
\end{equation}
where $R(\omega)$ is the greybody factor.
In the linear-dispersion regime ($k\xi \ll 1$), $R(\omega) \approx 1$ for
the dominant Hawking modes.
The Bose--Einstein occupation number then gives
\begin{equation}
  n(\omega) = \frac{1}{e^{\hbar\omega/\kB\TH} - 1},
\end{equation}
and the Fourier transform of the density-density correlator at the horizon
should exhibit a pole at $\omega = i\,\lambdaL$ in the complex frequency
plane.

\subsection{Protocol for measurement}

A possible measurement protocol:
\begin{enumerate}
  \item Create the analogue black hole as in \cite{steinhauer2019}.
  \item Measure $g^{(2)}(t, t+\tau)$ as a function of $\tau$ for phonons
    detected outside the sonic horizon.
  \item Fit the exponential tail: $g^{(2)}(\tau) - 1 = A\,e^{-\Gamma\tau}
    + B\,e^{-2\Gamma\tau} + \ldots$
  \item Compare $\Gamma$ with $\lambdaL = 2\pi\kB\TH/\hbar$ using the
    independently measured $\TH$.
\end{enumerate}
The ECI prediction is that $\Gamma = \lambdaL$ to within experimental
precision.
A significant departure (say $|\Gamma - \lambdaL|/\lambdaL > 10\%$) would
falsify the type~II$_\infty$ modular Lyapunov saturation picture.

%=======================================================
\section{Comparison with Steinhauer 2019 Data}
\label{sec:comparison}
%=======================================================

\subsection{What Steinhauer measured}

Steinhauer~et~al.\ (2019) \cite{steinhauer2019} measured:
\begin{itemize}
  \item The correlation spectrum $G^{(2)}(k, k')$ of Hawking-partner pairs
    in momentum space.
  \item The Planck-distributed occupation number $n(\omega)$ as a function
    of $\omega$, confirming thermality at $\TH = 0.35\,\mathrm{nK}$.
  \item The absence of firewall correlations (partner particles inside the
    horizon are only negative-energy partners, not Hawking particles).
\end{itemize}

\subsection{Overlap with ECI prediction}

The ECI prediction~\eqref{eq:prediction} overlaps with Steinhauer's
measurement in the following sense:
\begin{enumerate}
  \item \textbf{Compatible:} Steinhauer's $\TH = 0.35\,\mathrm{nK}$
    directly determines $\lambdaL = 2.88 \times 10^2\,\mathrm{s}^{-1}$.
    There is no contradiction.
  \item \textbf{Non-trivial extension:} Steinhauer did not measure the
    \emph{temporal} decay rate $\Gamma$ of $g^{(2)}(\tau)$.
    The prediction $\Gamma = \lambdaL$ is therefore an independent
    prediction of ECI, beyond what was measured.
  \item \textbf{Not yet a cross-check:} Since $\Gamma$ has not been
    measured, there is no existing data to compare against.
\end{enumerate}

The key implication: if a future measurement finds $\Gamma =
(2.88 \pm 0.5) \times 10^2\,\mathrm{s}^{-1}$ (consistent with
$\lambdaL$), this would constitute a non-trivial corroboration of the
ECI type~II$_\infty$ / modular Lyapunov saturation picture.
If $\Gamma \neq \lambdaL$, ECI is falsified on sonic-horizon
thermodynamics.

\subsection{Recent analogue gravity literature}

The community is actively developing new probes.
Chandran and Fischer (2026) \cite{chandran2026} show that logarithmic
entanglement negativity acquires a UV-finite volume term from Hawking
correlations --- this volume-law scaling is a potential experimental target.
The ECI type~II$_\infty$ framework predicts specific entanglement scaling
from the crossed-product trace; checking consistency with
\cite{chandran2026} is a required next step.

%=======================================================
\section{Open Questions and Signal-to-Noise Estimate}
\label{sec:open}
%=======================================================

\subsection{Signal-to-noise for $g^{(2)}(\tau)$ measurement}

The predicted decay rate $\Gamma \approx 288\,\mathrm{s}^{-1}$
corresponds to a decay time of $\sim 3.5\,\mathrm{ms}$.
Key experimental considerations:
\begin{itemize}
  \item \textbf{BEC lifetime:} Steinhauer's condensates survive for
    $\sim 100\,\mathrm{ms}$, sufficient to observe $\sim 30$ e-folding times.
  \item \textbf{Phonon coherence length:} Correlation measurements are
    performed in situ via density-density correlations in absorption
    imaging; temporal resolution $\sim 10$--$100\,\mu\mathrm{s}$ is achievable.
  \item \textbf{Signal-to-noise:} Each experimental run gives one
    shot-to-shot correlation; Steinhauer's 2019 paper uses $\sim 4600$
    shots. For temporal correlations, time-ordered pairs of shots are
    needed, requiring a modified experimental protocol.
  \item \textbf{Key bottleneck:} The characteristic frequency
    $\nu_{\mathrm{char}} = 7.3\,\mathrm{Hz}$ is very low; the phonon modes
    relevant for Krylov saturation have periods $\sim 137\,\mathrm{ms}$,
    comparable to the BEC lifetime.
    This is a marginal but potentially accessible regime.
\end{itemize}

\subsection{Trans-Planckian problem}

The BEC analogue evades the trans-Planckian problem \cite{jacobson1991}
by construction: the UV cutoff at $\xi$ prevents trans-$\xi$ modes from
contributing to Hawking radiation.
For the Krylov complexity computation, the relevant modes are $k \ll 1/\xi$
(infrared), so the result is robust against UV physics.
This is a major advantage of the BEC analogue over any thought experiment
in real gravity.

\subsection{Verch folium and BEC quasi-vacuum}

The BEC quasi-vacuum $|\Omega_{\mathrm{BdG}}\rangle$ is not in the
Verch folium of the conformal vacuum \cite{verch1994}: the Bogoliubov
transformation relating them is non-Hilbert-Schmidt at finite $\xi$.
However, the KMS property at $\TH$ holds for $k\xi \ll 1$, which suffices.
A rigorous operator-algebraic treatment of the BEC analogue in the spirit
of \cite{witten2022} is an open problem.

\subsection{Hartnoll--Yang BKL connection}

Hartnoll and Yang \cite{hartnoll2025} show that BKL dynamics near
cosmological singularities can be mapped to a conformal quantum mechanics
with modular invariance.
A potential connection to the BEC sonic horizon: the near-horizon region
has a BKL-like structure in the interior (supersonic region approaching
the inner horizon in Steinhauer's two-horizon setup).
This connection is speculative and requires detailed analysis.

\subsection{Contact with experimental groups}

The Steinhauer group (Technion) is the primary contact.
Their 2024--2026 arXiv output has shifted to pure GR topics; no
BEC-sonic-horizon saturation-envelope paper exists on arXiv as of
2026-05-03 (see \texttt{notes.md} for full search results).
The Munoz de Nova group (Universidad Complutense, Madrid) continues active
work on BEC analogue gravity and would be a complementary contact.

%=======================================================
\section{Retraction Flag: ECI v6.0.10 ``8.29\% Saturation''}
\label{sec:retraction}
%=======================================================

\textbf{Flag for ECI~v6.0.30 retraction.}

The ECI~v6.0.10 claim:
\begin{quote}
``$\rho \approx 8.29\%$ at leading order, consistent across three
Steinhauer datasets within $1$--$10\%$ correction envelope.''
\end{quote}
is \textbf{unverifiable} based on the following evidence:

\begin{enumerate}
  \item No Steinhauer paper (arXiv:1510.00621, 1809.00913, 1910.09363)
    reports a dimensionless ratio $\rho$ or a ``saturation envelope.''
  \item The arXiv API search for Steinhauer-group 2024--2026 papers on
    BEC saturation returned \textbf{zero results} for any such claim.
  \item The ECI framework does not define what observable $\rho$ measures
    in Steinhauer's experimental units.
  \item This claim has the hallmarks of an LLM-hallucinated benchmark
    (project history records 28+ such incidents).
\end{enumerate}

\textbf{Action:} Remove this claim from all ECI documents and submitted
papers.
Replace with the honest ECI prediction \eqref{eq:prediction}, which is
new and testable but has \emph{not yet been checked against data}.

%=======================================================
\section*{References}
%=======================================================

\begin{thebibliography}{99}

\bibitem{unruh1981}
W.~G.~Unruh,
\emph{Experimental black-hole evaporation?}
Phys.\ Rev.\ Lett.\ \textbf{46}, 1351 (1981).
[No arXiv; pre-arXiv era.]

\bibitem{steinhauer2016}
J.~Steinhauer,
\emph{Observation of quantum Hawking radiation and its entanglement in an analogue black hole},
Nat.\ Phys.\ \textbf{12}, 959 (2016).
arXiv:1510.00621 [gr-qc].

\bibitem{steinhauer2019}
J.~R.~Mu{\~n}oz de Nova, K.~Golubkov, V.~I.~Kolobov, J.~Steinhauer,
\emph{Observation of thermal Hawking radiation at the Hawking temperature in an analogue black hole},
Nature \textbf{569}, 688 (2019).
arXiv:1809.00913 [gr-qc, cond-mat.quant-gas].

\bibitem{kolobov2021}
V.~I.~Kolobov, K.~Golubkov, J.~R.~Mu{\~n}oz de Nova, J.~Steinhauer,
\emph{Observation of stationary spontaneous Hawking radiation and the time evolution of an analogue black hole},
Nat.\ Phys.\ \textbf{17}, 362 (2021).
arXiv:1910.09363 [gr-qc].

\bibitem{jacobson1991}
T.~Jacobson,
\emph{Black-hole evaporation and ultrashort distances},
Phys.\ Rev.\ D~\textbf{44}, 1731 (1991).
[Pre-arXiv; no arXiv ID verified.]

\bibitem{witten2022}
E.~Witten,
\emph{Gravity and the crossed product},
JHEP \textbf{2022}, 8 (2022).
arXiv:2112.12828 [hep-th].

\bibitem{mss2016}
J.~Maldacena, S.~H.~Shenker, D.~Stanford,
\emph{A bound on chaos},
JHEP \textbf{2016}, 106 (2016).
arXiv:1503.01409 [hep-th].

\bibitem{caputa2023}
P.~Caputa, J.~M.~Mag\'an, D.~Patramanis, E.~Tonni,
\emph{Krylov complexity of modular Hamiltonian evolution},
(2023).
arXiv:2306.14732 [hep-th].

\bibitem{hartnoll2025}
S.~A.~Hartnoll, M.~Yang,
\emph{The conformal primon gas at the end of time},
(2025).
arXiv:2502.02661 [hep-th].

\bibitem{delhom2025}
A.~Delhom, L.~Giacomelli,
\emph{Analogue gravity with Bose-Einstein condensates},
(2025).
arXiv:2512.14209 [gr-qc, cond-mat.quant-gas, quant-ph].

\bibitem{vardian2026}
N.~Vardian,
\emph{Modular Krylov complexity as a boundary probe of area operator and entanglement islands},
(2026).
arXiv:2602.02675 [hep-th].

\bibitem{chandran2026}
S.~M.~Chandran, U.~R.~Fischer,
\emph{Emergence of volume-law scaling for entanglement negativity from the Hawking radiation of analogue black holes},
(2026).
arXiv:2604.02075 [gr-qc, cond-mat.quant-gas, quant-ph].

\bibitem{anderson2024}
P.~R.~Anderson, R.~Balbinot, R.~A.~Dudley, A.~Fabbri,
\emph{Looking for traces of Hawking radiation in correlation functions of BEC acoustic black holes},
Comptes Rendus Physique (2025).
arXiv:2410.02700 [gr-qc].

\bibitem{munozdenova2024}
J.~R.~Mu{\~n}oz de Nova, P.~Fern\'andez Palacios, P.~Alc\'azar Guerrero, I.~Zapata, F.~Sols,
\emph{Resonant analogue configurations in atomic condensates},
(2024).
arXiv:2406.10027 [cond-mat.quant-gas].

\bibitem{verch1994}
R.~Verch,
\emph{Nuclearity, split property, and duality for the Klein-Gordon field in curved spacetime},
Lett.\ Math.\ Phys.\ \textbf{29}, 297 (1994).
[Verch folium reference; no arXiv ID verified.]

\bibitem{eci_krylov_diameter}
ECI Collaboration,
\emph{Krylov diameter and generalized entropy: Theorem 4},
internal preprint v6.0.15 (2025).

\end{thebibliography}

\end{document}
